% ---------------------------------------------------------------------
%  CHƯƠNG 4 — THIẾT KẾ HỆ THỐNG VÀ THUẬT TOÁN
% ---------------------------------------------------------------------
\chapter{Thiết kế hệ thống và Thuật toán}
\label{ch:thietke}

\begin{muctieu}
Chương này trình bày chi tiết thiết kế kỹ thuật của hệ thống FoodAI, bao gồm kiến trúc 3 tầng (Three-Tier Architecture), mô hình cơ sở dữ liệu quan hệ chuẩn 3NF (9 bảng), cùng thiết kế giải thuật gợi ý món ăn theo kiến trúc Database-First RAG và thuật toán kiểm soát hình ảnh Cookpad chống sai lệch.
\end{muctieu}

\section{Thiết kế kiến trúc hệ thống}

Hệ thống được thiết kế theo mô hình kiến trúc ba tầng \en{(Three-Tier Architecture)} hiện đại, đảm bảo tính mô-đun hóa, bảo mật và khả năng mở rộng dễ dàng (Hình~\ref{fig:kientruc_hethong}).

\begin{figure}[H]
\centering
\resizebox{\textwidth}{!}{%
\begin{tikzpicture}[node distance=10mm and 12mm]
  % Tầng Client
  \node[startstop, fill=lightnavy, draw=daunavy] (client) {\textbf{TẦNG GIAO DIỆN (CLIENT LAYER)}\\React 18 + TypeScript + Vite + Tailwind CSS\\(User Dashboard \& Admin Dashboard)};

  % Tầng Server
  \node[process, fill=white, draw=daumaroon, below=12mm of client] (server) {\textbf{TẦNG ỨNG DỤNG (SERVER LAYER)}\\Node.js + Express RESTful API\\Auth Middleware (JWT) $\cdot$ Rate Limiter $\cdot$ CORS};

  % Khối dịch vụ bên trong Server
  \node[clsbox, below left=10mm and -5mm of server] (ai_svc) {Gemini Service\\(AI Chef \& Nutritionist)};
  \node[clsbox, below=10mm of server] (recipe_svc) {Recipe Service\\(Heuristic Matching)};
  \node[clsbox, below right=10mm and -5mm of server] (img_svc) {Image Service\\(Multi-source Retrieval)};

  % Tầng Dữ liệu & AI bên ngoài
  \node[objbox, below=14mm of recipe_svc] (db) {\textbf{CƠ SỞ DỮ LIỆU}\\MySQL 8.0 / SQLite\\(Users, Recipes, Cookpad Images, Favorites, History)};

  \node[membox, left=10mm of db] (gemini_ext) {\textbf{Google Gemini API}\\(LLM Engine)};
  \node[membox, right=10mm of db] (pexels_ext) {\textbf{Pexels \& Cookpad}\\(Image Sources)};

  % Kết nối
  \draw[flowarr] (client) -- node[right, font=\tiny]{HTTPS / JSON} (server);
  \draw[flowarr] (server) -- (ai_svc);
  \draw[flowarr] (server) -- (recipe_svc);
  \draw[flowarr] (server) -- (img_svc);
  \draw[flowarr] (ai_svc) -- (gemini_ext);
  \draw[flowarr] (recipe_svc) -- (db);
  \draw[flowarr] (img_svc) -- (pexels_ext);
  \draw[flowarr] (img_svc) -- (db);
\end{tikzpicture}}
\caption{Sơ đồ kiến trúc phân tầng tổng thể của hệ thống FoodAI.}
\label{fig:kientruc_hethong}
\end{figure}

\noindent\textbf{Mô tả:} Hình~\ref{fig:kientruc_hethong} mô tả kiến trúc phân tầng tổng thể của hệ thống FoodAI theo mô hình 3 tầng (Three-Tier Architecture), bao gồm: (1) Tầng giao diện Client Layer (React 18, TypeScript, Vite, Tailwind CSS); (2) Tầng ứng dụng Server Layer (Node.js, Express RESTful API với các khối Auth Middleware, Rate Limiter, Gemini Service, Recipe Service, Image Service); và (3) Tầng dữ liệu \& dịch vụ bên ngoài (CSDL MySQL 8.0, Google Gemini API, Pexels API và kho ảnh Cookpad).

\noindent\textbf{Diễn giải:} Sơ đồ thể hiện tính phân tách trách nhiệm (Separation of Concerns) cao giữa các tầng. Mọi yêu cầu từ Client đều phải thông qua giao thức HTTPS có kiểm soát bởi Middleware xác thực và giới hạn tần suất (Rate Limiter) trước khi chạm tới các dịch vụ nghiệp vụ. Việc bóc tách các Service độc lập (Gemini Service cho AI, Recipe Service cho giải thuật so khớp, Image Service cho truy xuất ảnh) giúp mã nguồn đạt độ gắn kết cao, giảm thiểu phụ thuộc lỏng lẻo, tạo điều kiện thuận lợi cho việc kiểm thử đơn vị, bảo trì và mở rộng tính năng mà không gây ảnh hưởng tới toàn bộ hệ thống.

\section{Thiết kế cơ sở dữ liệu}

Cơ sở dữ liệu quan hệ của hệ thống được thiết kế chuẩn hóa để lưu trữ thông tin người dùng, công thức ẩm thực, hình ảnh cào từ Cookpad và dữ liệu tương tác (Bảng~\ref{tab:schema_users} và Bảng~\ref{tab:schema_recipes}).

\begin{table}[H]
\centering
\caption{Lược đồ cấu trúc bảng người dùng (\code{users}).}
\label{tab:schema_users}
\begin{tabularx}{\textwidth}{@{}>{\bfseries\color{daunavy}}p{2.6cm} >{\color{daumaroon}}p{2.8cm} >{\raggedright\arraybackslash}X@{}}
\toprule
\rowcolor{gray100}\bfseries\color{daunavy}Tên cột & \bfseries\color{daumaroon}Kiểu dữ liệu & \bfseries\color{daunavy}Ràng buộc \& Ý nghĩa \\
\midrule
id          & VARCHAR(36)  & Khóa chính (UUID v4) \\
username    & VARCHAR(50)  & Tên đăng nhập duy nhất (UNIQUE, NOT NULL) \\
email       & VARCHAR(100) & Địa chỉ email duy nhất (UNIQUE, NOT NULL) \\
password    & VARCHAR(255) & Mật khẩu đã băm (Bcrypt hash) \\
full\_name  & VARCHAR(100) & Họ và tên hiển thị \\
role        & VARCHAR(10)  & Vai trò tài khoản (\code{'user'} hoặc \code{'admin'}) \\
is\_active  & TINYINT(1)   & Trạng thái kích hoạt tài khoản (1: Hoạt động, 0: Khóa) \\
created\_at & DATETIME     & Thời điểm tạo tài khoản \\
\bottomrule
\end{tabularx}
\end{table}

\noindent\textbf{Mô tả:} Bảng~\ref{tab:schema_users} mô tả chi tiết lược đồ cấu trúc của bảng \code{users} gồm 8 trường dữ liệu: id (khóa chính UUID), username (tên đăng nhập duy nhất), email (địa chỉ thư điện tử duy nhất), password (chuỗi băm BCrypt 255 ký tự), full\_name (họ tên hiển thị), role (vai trò người dùng), is\_active (cờ trạng thái kích hoạt) và created\_at (mốc thời gian tạo tài khoản).

\noindent\textbf{Diễn giải:} Thiết kế bảng đáp ứng chuẩn dạng chuẩn 3 (3NF) trong cơ sở dữ liệu quan hệ. Ràng buộc UNIQUE trên hai trường username và email ngăn chặn tuyệt đối tình trạng đăng ký trùng lặp tài khoản. Việc lưu trữ mật khẩu bằng chuỗi băm BCrypt 255 ký tự đảm bảo an toàn tối đa theo các tiêu chuẩn an ninh mạng hiện đại. Trường cờ is\_active cho phép quản trị viên có thể vô hiệu hóa quyền truy cập của tài khoản vi phạm chính sách một cách linh hoạt mà vẫn bảo toàn toàn bộ dữ liệu lịch sử tìm kiếm và món ăn yêu thích liên kết.

\begin{table}[H]
\centering
\caption{Lược đồ cấu trúc bảng món ăn nội bộ (\code{recipes}).}
\label{tab:schema_recipes}
\begin{tabularx}{\textwidth}{@{}>{\bfseries\color{daunavy}}p{2.8cm} >{\color{daumaroon}}p{2.8cm} >{\raggedright\arraybackslash}X@{}}
\toprule
\rowcolor{gray100}\bfseries\color{daunavy}Tên cột & \bfseries\color{daumaroon}Kiểu dữ liệu & \bfseries\color{daunavy}Ràng buộc \& Ý nghĩa \\
\midrule
id            & VARCHAR(36)  & Khóa chính món ăn \\
name          & VARCHAR(255) & Tên món ăn tiếng Việt (INDEX) \\
description   & TEXT         & Mô tả vắn tắt về hương vị \\
image\_url    & VARCHAR(1000)& Đường dẫn ảnh món ăn thực tế \\
cooking\_time & INT          & Thời gian chế biến (phút) \\
difficulty    & VARCHAR(20)  & Độ khó (\code{'Dễ'}, \code{'Trung bình'}, \code{'Khó'}) \\
category      & VARCHAR(50)  & Danh mục (\code{'Món xào'}, \code{'Súp/Canh'}, \dots) \\
author\_name  & VARCHAR(100) & Tác giả chia sẻ công thức trên Cookpad \\
\bottomrule
\end{tabularx}
\end{table}

\noindent\textbf{Mô tả:} Bảng~\ref{tab:schema_recipes} mô tả cấu trúc bảng lưu trữ món ăn \code{recipes} với 8 thuộc tính nòng cốt: id (khóa chính), name (tên món ăn tiếng Việt có đánh chỉ mục INDEX), description (mô tả hương vị), image\_url (đường dẫn hình ảnh CDN), cooking\_time (thời gian nấu ước tính tính theo phút), difficulty (độ khó chế biến), category (danh mục phân loại) và author\_name (tác giả chia sẻ công thức).

\noindent\textbf{Diễn giải:} Bảng \code{recipes} đóng vai trò là cơ sở tri thức ẩm thực nội bộ (Internal Knowledge Base) của hệ thống. Việc đánh chỉ mục INDEX trên trường tên món ăn \code{name} giúp tối ưu hóa đáng kể tốc độ tìm kiếm văn bản và so khớp Heuristic khi hệ thống hoạt động ở chế độ Fallback, giảm thời gian phản hồi từ CSDL xuống dưới mức 50ms. Các trường định lượng như thời gian nấu và mức độ khó cung cấp siêu dữ liệu đầy đủ để giao diện người dùng hiển thị thẻ món ăn trực quan và hỗ trợ bộ lọc đa chiều.

Ngoài ra, hệ thống lưu trữ ảnh trích xuất trong bảng \code{cookpad\_dish\_images}, lưu trữ lịch sử người dùng trong \code{search\_history} và lưu công thức yêu thích dạng JSON chi tiết trong \code{favorite\_recipes}.

\section{Thiết kế thuật toán gợi ý và Xử lý ảnh thông minh}

\subsection{Luồng thực thi gợi ý món ăn thông minh}

Quy trình gợi ý món ăn kết hợp AI và Fallback được mô tả chi tiết tại Hình~\ref{fig:flow_suggest}.

\begin{figure}[H]
\centering
\resizebox{0.82\textwidth}{!}{%
\begin{tikzpicture}[node distance=10mm]
  \node[startstop] (start) {Bắt đầu: Tiếp nhận danh sách nguyên liệu $U$};
  \node[decision, below=8mm of start] (check_key) {Cấu hình Gemini API hợp lệ?};
  \node[process, below left=10mm and 2mm of check_key] (call_ai) {Gọi Google Gemini AI\\với Prompt ràng buộc};
  \node[decision, below=8mm of call_ai] (check_ai) {AI phản hồi \&\\Parse JSON thành công?};
  \node[process, below=8mm of check_ai] (post_filter) {Hậu kiểm: Khử trùng lặp\\và lọc vi phạm thịt cá};
  \node[process, right=20mm of check_ai] (fallback_db) {Kích hoạt chế độ Fallback:\\Tra cứu kho Cookpad MySQL};
  \node[process, below=8mm of post_filter] (attach_img) {Truy xuất \& Gán hình ảnh độc nhất\\từ Cookpad / Pexels / HD Bank};
  \node[startstop, below=8mm of attach_img] (finish) {Trả về danh sách 4--6 món ăn kèm\\công thức, dinh dưỡng và ảnh cho người dùng};

  \draw[flowarr] (start) -- (check_key);
  \draw[flowarr] (check_key) -| node[above, font=\footnotesize]{Có} (call_ai);
  \draw[flowarr] (check_key) -| node[above, font=\footnotesize]{Không} (fallback_db);
  \draw[flowarr] (call_ai) -- (check_ai);
  \draw[flowarr] (check_ai) -- node[right, font=\footnotesize]{Thành công} (post_filter);
  \draw[flowarr] (check_ai) -- node[above, font=\footnotesize]{Lỗi / Timeout} (fallback_db);
  \draw[flowarr] (fallback_db) |- (attach_img);
  \draw[flowarr] (post_filter) -- (attach_img);
  \draw[flowarr] (attach_img) -- (finish);
\end{tikzpicture}}
\caption{Sơ đồ khối luồng xử lý gợi ý món ăn thông minh (Hybrid Flowchart).}
\label{fig:flow_suggest}
\end{figure}

\noindent\textbf{Mô tả:} Hình~\ref{fig:flow_suggest} mô tả sơ đồ lưu đồ (Flowchart) tiến trình xử lý gợi ý món ăn, thể hiện các khối quyết định rẽ nhánh logic từ lúc tiếp nhận danh sách nguyên liệu đầu vào, kiểm tra điều kiện khả dụng của API Key, gọi mô hình Gemini AI, kiểm tra tính hợp lệ của gói tin JSON, tiến hành lọc khử vi phạm nguyên liệu, cơ chế kích hoạt nhánh Fallback cơ sở dữ liệu khi phát sinh lỗi, đến công đoạn gắn hình ảnh và trả về kết quả cuối cùng.

\noindent\textbf{Diễn giải:} Sơ đồ khối giải thích nguyên nhân hệ thống duy trì được tính ổn định và tính chính xác cao. Nhánh xử lý AI được kiểm soát chặt chẽ bởi hai chốt chặn: khối kiểm tra cú pháp JSON đảm bảo cấu trúc dữ liệu không bị gãy vỡ, và khối hậu kiểm ngữ nghĩa loại bỏ ngay các món ăn tự bịa thịt cá vi phạm ràng buộc nguyên liệu đầu vào. Nhánh rẽ Fallback sang kho Cookpad MySQL giải thích vì sao khi mất kết nối mạng hay hết quota Gemini, hệ thống vẫn trả về được các món ăn truyền thống chuẩn vị và có sẵn ảnh chất lượng cao mà không hiển thị màn hình trắng hay báo lỗi khó chịu cho người dùng.

\subsection{Giải thuật gợi ý món ăn theo kiến trúc Database-First RAG (Mã giả)}

Giải thuật gợi ý món ăn áp dụng kỹ thuật RAG định hướng CSDL (\en{Database-First Retrieval-Augmented Generation})~\cite{Lewis2020RAG} kết hợp làm giàu dữ liệu tự động bằng Google Gemini AI được đặc tả chi tiết tại Giải thuật~\ref{alg:hybrid_suggest}.

\begin{algorithm}[H]
\caption{Gợi ý món ăn thông minh Database-First RAG và AI Enrichment}
\label{alg:hybrid_suggest}
\begin{algorithmic}[1]
\Require Danh sách nguyên liệu người dùng $U$, số lượng món cần lấy $K = 8$
\Ensure Danh sách $R$ gồm các món ăn có tên, hình ảnh chuẩn MySQL, công thức và dinh dưỡng do AI phân tích
\State Khởi tạo danh sách kết quả $R \gets \emptyset$
\State \Comment{Bước 1: Truy vấn CSDL lọc món ăn thực tế và ảnh Cookpad theo logic nguyên liệu}
\State $Cand \gets$ \Call{SuggestRecipesFromDatabase}{$U$, $\max(2K, 10)$}
\If{\Call{IsGeminiEnabled}{} = \textbf{True}}
    \State \textbf{thử thực hiện (try):}
        \State \quad \Comment{Bước 2: Gửi danh sách món từ CSDL cho Gemini AI phân tích chuyên sâu}
        \State \quad $TargetDishes \gets Cand[0 \dots K-1]$
        \State \quad $P \gets$ \Call{BuildEnrichmentPrompt}{$U$, $TargetDishes$}
        \State \quad $RawResponse \gets$ \Call{GeminiGenerate}{$P$}
        \State \quad $ParsedDishes \gets$ \Call{ParseAndSanitizeJSON}{$RawResponse$}
        \State \quad $EnrichedRecipes \gets$ \Call{AlignWithDatabaseCandidates}{$ParsedDishes$, $TargetDishes$}
        \State \quad \Comment{Bước 3: Tự động đồng bộ và lưu công thức, dinh dưỡng AI về lại MySQL}
        \State \quad \Call{SaveAiRecipesToDatabase}{$EnrichedRecipes$}
        \State \quad \Return $EnrichedRecipes$ \Comment{Thành công ở tầng AI Enrichment}
    \State \textbf{bắt lỗi (catch):}
        \State \quad Ghi nhật ký lỗi và chuyển sang chế độ CSDL thuần (Fallback)
\EndIf
\State \Comment{Bước 4 (Dự phòng): Trả về danh sách ứng viên trực tiếp từ kho Cookpad MySQL}
\State $R_{db} \gets Cand[0 \dots K-1]$
\State \Return $R_{db}$
\end{algorithmic}
\end{algorithm}

\noindent\textbf{Mô tả:} Giải thuật~\ref{alg:hybrid_suggest} mô tả mã giả của quy trình gợi ý món ăn kết hợp Database-First Retrieval-Augmented Generation (RAG) gồm 4 bước: (1) Truy vấn danh sách ứng viên món ăn và ảnh thật từ kho MySQL; (2) Truyền ngữ cảnh ứng viên cho Gemini AI để phân tích dinh dưỡng và hướng dẫn nấu; (3) Đồng bộ kết quả AI ngược về CSDL; (4) Nhánh dự phòng trả về dữ liệu nội bộ khi AI gặp lỗi.

\noindent\textbf{Diễn giải:} Giải thuật này lý giải bản chất logic vì sao hệ thống loại bỏ được nhược điểm lớn nhất của mô hình ngôn ngữ lớn là gợi ý các món ăn không có trong thực tế hoặc bịa đặt ảnh minh họa. Bằng cách định hướng lấy ứng viên từ CSDL trước (Database-First), các món ăn và hình ảnh luôn được neo chặt vào dữ liệu ẩm thực thực tế từ Cookpad. AI chỉ đảm nhận vai trò phân tích công thức, mẹo nấu và tính toán dinh dưỡng (Enrichment). Việc tự động lưu ngược kết quả AI vào MySQL giúp hệ thống dần tích lũy tri thức, biến kho dữ liệu ngày càng phong phú sau mỗi lượt truy vấn.
